\documentclass[11pt,a4paper]{article}
\usepackage[T1]{fontenc}
\usepackage{lmodern,amsmath,amssymb,amsthm,microtype,booktabs}
\usepackage[margin=27mm]{geometry}
\usepackage[colorlinks=true,linkcolor=blue,urlcolor=blue]{hyperref}
\setlength{\emergencystretch}{2em}
\newtheorem{proposition}{Proposition}
\newcommand{\Q}{\mathbb Q}\DeclareMathOperator{\Tr}{Tr}
\title{Failure of automatic preparation for raw abscissa factors}
\author{}\date{17 November 2026}
\begin{document}\maketitle
\begin{abstract}
At an exact multiplication tuple and an actual physical slope, replacing
four balanced abscissa factors by either $1+C_i/u$ or $1+C_i^2/u$
creates negative odd valuation in the pole input. The two comparison
tests remain balanced. All pre-pole point, contact and detector guards
hold, and the real pole input lies strictly between zero and one at
every point character. The defect occurs in both conjugate slots at
$76777$. Projective multiplication modulo $76777^2$ and a separate
affine computation certify it. This disproves automatic preparation
for these raw replacements; it does not exclude a different choice
of witnesses or a new completeness argument.
\end{abstract}

\section*{The two precisely specified replacements}
Retain the derivative detector's curve and non-torsion point
\[
 K=\Q(\sqrt2),\quad E:y^2=x^3-d^2x,\quad d=3-\sqrt2,
 \quad P=\left(\frac{57-28\sqrt2}{8},\frac{-660+417\sqrt2}{32}\right).
\]
The arithmetic input $2E(K)=\mathbb ZP$ and the point charts are those
of \cite{Derivative,FullAux}. Put $f=1/x$, $\zeta=-x/y$, and retain
\[
 r=\frac{\zeta(U)}{\zeta(H)},\qquad
 r'=\overline{\frac{\zeta(U-A)}{\zeta(H-P)}},
\]
\[
 (C_1,\ldots,C_6)
 =\bigl(f(H),f(U),\overline{f(H-P)},\overline{f(U-A)},r-r',r-X\bigr).
\]
For $e\in\{1,2\}$ define the candidate raw detector
\begin{equation}\label{eq:raw}
 W_e=\prod_{i=1}^4\left(1+\frac{C_i^e}{u}\right)
          \prod_{i=5}^6\left(1+\frac{D_i^2}{u}\right),
 \qquad D_i=\frac{C_i^2}{1+uC_i^2}.
\end{equation}
Thus the exponent-one variant uses the even abscissa valuations directly;
the exponent-two variant merely omits the first four balancing denominators.
Both retain the original balanced comparison tests. Write
\[
 V_e=\Tr(W_e^{-1}),\quad w=\frac{u}{(1+u)^2},\quad
 b_e=\frac{4w}{V_e^2},\quad\eta_e=7b_e,
 \qquad R_e=V_e^2+w(2-V_e).
\]

\section*{An exact actual tuple}
\begin{proposition}
Take
\[
 X=n=2,\qquad B=3,\qquad A=2P,\quad H=10336P,\quad U=20672P=2H.
\]
The root trace and the physical data are
\[
 T=\Tr x(2P)=\frac{8226963211}{1142622768},\qquad
 M=7^4+T^4,\qquad a=M^4,\qquad z=670434583579295,
\]
\[
 F=\prod_{c\in\{-10,-8,-5,1,2,3,4,6,7\}}((a+c)^2-z^2),
 \qquad u=-18F/F_a.
\]
Here $F_a$ is differentiated before specializing $a$.
All pre-pole point, difference, contact and detector denominator guards
are nonzero at every point-algebra character. In both real embeddings,
and at every real point character, one has
\[
 z>a+100,\qquad u>100,\qquad 0<\eta_e<\frac{847}{10000}<1.
\]
Nevertheless, at $p=76777$,
\[
                 v_p(b_1)=-1,\qquad v_p(b_2)=-5,
                 \qquad v_p(R_1)=v_p(R_2)=1.
\]
Thus neither actual $R_e$ is a rational square, and neither $\eta_e$
has square reduced denominator. The intrinsic pole equation has no
rational solution for either input.
\end{proposition}

\begin{proof}
Every point or difference used here is a nonzero multiple of $P$.
In particular the difference indices are $10335$ and $20670$.
Non-torsion excludes zero abscissae and ordinates. The affine addition
denominators are nonzero because $10336\ne\pm1$ and
$20672\ne\pm2$. The same statements hold after independently
changing the ordinate signs in either real slot: the resulting
difference indices are among $\pm10336-1$ and
$\pm20672\mp2$, none of which is zero. All these points lie in
$2E(K)$.

For completeness of the real guard argument, every nonzero double on
$y^2=x^3-d^2x$ has abscissa at least $d$ in its real embedding.
This follows, for example, from
\[
 x(2Q)-d=
 \left(\frac{x(Q)^2-2d x(Q)-d^2}{2y(Q)}\right)^2\ge0.
\]
Here $d$ is positive in both embeddings, and $3-\sqrt2>1$.
Equality would make $2Q$ a point of order two, which cannot occur
for a nonzero multiple of $P$. Therefore each of the first four
$C_i$ lies strictly between zero and one, at every real point
character. The ordinate roots of the point algebra at these data
are real in both slots. Thus these sign characters exhaust all its
characters for purposes of the nonzero norm guards.

Exact rational comparison gives $z>a+100$. Every contact factor is
negative, so $F<0$, $F_a>0$, and $u>0$. More precisely,
\[
 -\frac{F_a}{F}
  =\sum_c\frac{2(a+c)}{z^2-(a+c)^2}
       <\frac9{z-a-7},\qquad u>2(z-a-7)>100.
\]
Hence $M,a^2-z^2,F,F_a,u,1+u$ are nonzero. Both comparison
denominators are positive and $0\le D_i<1/u$ for $i=5,6$.
Each real slot of \eqref{eq:raw} consequently satisfies
\[
 1<W_e<\left(\frac{101}{100}\right)^4
              \left(1+\frac1{10^6}\right)^2<\frac{11}{10}.
\]
Thus $20/11<V_e<2$, $0<w<1/100$, and
$0<\eta_e<7\cdot4(1/100)/(20/11)^2=847/10000$.
Every product, trace and comparison-denominator norm is nonzero.
These arguments establish all pre-pole guards, including at the other
point characters; no pole ordinate is asserted to exist.

We now give the exact local certificate. The integer $p=76777$ is
prime and exceeds the source's $76176$ cutoff. Put
$p^2=5894707729$. The two roots of two modulo $p^2$ used for the
slots are
\[
                    3367298917,\qquad2527408812.
\]
They reduce to $13251$ and its negative modulo $p$. Independent
elliptic multiplication gives the following data for $Q=H-P=10335P$:
\[
\begin{array}{c|r|r|r}
 \sqrt2\pmod {p^2}&x(Q)\pmod {p^2}& y(Q)/p\pmod p&x(Q)/p^2\pmod p\\
 \hline
 3367298917&0&57786&64192\\
 2527408812&0&63990&61852
\end{array}
\]
The two middle nonzero residues show $v(y(Q))=1$. Since $d^2$ is
a unit, the curve equation gives $v(x(Q))=2$ and
$x(Q)/p^2\equiv-(y(Q)/p)^2/d^2\pmod p$, yielding the last column.
In each residue slot $H=P+(0,0)$, $U=A$, and $U-A$ reduces to
the identity. The abscissae of $H,U$ are units. Duplication near
$(0,0)$, or the formula
\[
            f(2Q)=\frac{4y(Q)^2}{(x(Q)^2+d^2)^2},
\]
shows $v(f(U-A))=2$. Thus the four abscissa-test valuations are
$(0,0,-2,2)$ in both slots.

Exact evaluation of the physical contact polynomial gives
\[
 z\equiv21480\pmod p,\qquad F\equiv44843\pmod p,
 \qquad F_a\equiv169753947\pmod {p^2}.
\]
Here $169753947=2211p$ with $2211\not\equiv0\pmod p$.
Consequently
\[
                  v_p(u)=-1,\qquad pu\equiv66932\pmod p.
\]
The two comparison factors in \eqref{eq:raw} are units congruent
to one: their old balanced valuation identity gives
$v(D_i^2/u)\ge3$ independently of possible poles in the comparisons.
Of the four raw factors, exactly the factor for $C_3$ has negative
valuation, namely $1-2e$; the other three are one modulo $p$.
Therefore $v(W_e^{-1})=2e-1$ in each slot.

Let $\xi_1=64192$, $\xi_2=61852$. Since $C_3$ is conjugated,
the two reciprocal-product leading coefficients are interchanged;
their sum is unchanged. They give
\[
 \frac{V_e}{p^{2e-1}}
       \equiv66932(\xi_1^e+\xi_2^e)\pmod p
       =\begin{cases}43471,&e=1,\\72176,&e=2.\end{cases}
\]
Both residues are nonzero. Hence $v_p(V_e)=2e-1$.
One also has $v_p(w)=1$, so
$v_p(b_e)=1-2(2e-1)=3-4e$.
Finally $2-V_e$ is a unit, while $v_p(V_e^2)\ge2$;
therefore $v_p(R_e)=1$.
The intrinsic pole implication excludes a negative odd valuation of
$\eta_e$; at $p\ne7$ its valuation equals that of $b_e$.
This proves the asserted failures.
\end{proof}

\section*{Scope and independent verification}
The example is an actual multiplication tuple and uses the actual
physical derivative quotient, not an independently chosen trace or
weight. It satisfies the full real inequalities and the pre-pole
nonvanishing guards. It is not asserted to satisfy the selected-prime
factorization or all fixed finite-place preparation conditions of a
completeness proof. The fixed shift $B=3$ is an algebraic fixture, not
a certified sufficient defining shift. The conclusion is that the
removed balancing factors do not prepare the new denominator primes
automatically. A different choice of auxiliary indices, stronger
prime-order avoidance, or another complete source is not excluded.

The exact specification and certificates are in
\nolinkurl{negative_results/computation/pole_parameters/raw_abscissa_v75}.
The producer uses projective multiplication modulo $p^2$, verifies the
curve equation and primitive reduction at every step, and rejects a
nonregular addition formula. The independent checker imports no
producer module and instead performs affine multiplication modulo $p^2$,
checking that every denominator it inverts is a unit. It reconstructs
$\Tr x(2P)$ from the original curve and evaluates $F_a$ by a separate
product-derivative sum. The large rational coordinates of $10336P$
are not needed: their group-multiple specification and good reduction
are exact.

A third checker evaluates the division-polynomial recurrences of
\cite[Sections 5.5--5.6]{Division} directly modulo $p^3$, importing
neither group-law implementation. It verifies the curve equation,
all inverted denominators, and a separate doubling control. The
literal values of $x(10335P)/p^2$ modulo $p$ are $64192$ and $61852$,
independently confirming the two decisive abscissa valuations.

\begin{thebibliography}{9}
\bibitem{Derivative}\emph{An eleven-variable construction under rational $abc$},
\nolinkurl{papers/conditional_few_variables/derivative_detector11_20261117_v63.tex}.
\bibitem{FullAux}\emph{Full auxiliary elliptic points in the eleven-variable source},
\nolinkurl{papers/conditional_few_variables/full_auxiliary11_20261117_v71.tex}.
\bibitem{Even}\emph{Even abscissa valuations and a smaller eleven-variable degree bound},
\nolinkurl{papers/conditional_few_variables/even_valuation_detector_20261117_v75.tex}.
\bibitem{Division}A. Sutherland, \emph{18.783 Elliptic Curves,
Lecture 5: Isogeny kernels and division polynomials}, MIT, Fall 2025,
Sections 5.5--5.6 and Theorem 5.21; checked 17 November 2026.
\url{https://math.mit.edu/classes/18.783/2025/LectureNotes5.pdf}.
\end{thebibliography}
\end{document}
